\section{从 sequential data 到 sequential decision making}

在进入 RL 公式之前，先理解讲者为什么从语言、视觉、蛋白质和代码生成讲起。这些领域的共同经验不是“Transformer 无所不能”，而是大规模监督或自监督序列建模通常拥有相对平滑的 loss、成熟的正则化手段和可复用的训练基础设施。讲者真正想迁移到决策问题中的，是这种训练接口与规模化路径。

\subsection{Transformer 成功的共同抽象}

语言、图像 patch、蛋白质残基和代码 token 看似不同，但都可以被表示为序列，并通过 masked prediction 或 autoregressive prediction 学习条件分布。读下面的 slide 时，先不要比较各模型的具体参数量；应先观察它们如何把领域对象变成 token，以及如何用同一类 next-token machinery 建模。

\lecturefigure{slide-02-transformers-for-sequential-data.jpg}{Transformer 已被用于语言、视觉、蛋白质与代码等多类 sequential data。}{Stanford Online 官方视频 01:40。}

\begin{knowledgebox}{读图：共同点不是输入模态，而是训练接口}
图中的 GPT-3、ViT、AlphaFold 与代码模型并不共享相同的输入语义，却都把复杂对象编码为可组合的表示，并通过大批量梯度优化学习条件关系。讲者借这组例子提出一个工程问题：如果决策轨迹也能转成序列，是否可以复用更稳定的训练目标、优化器和 scaling recipe？这是一项研究假设，不是由这些成功案例自动推出的定理。
\end{knowledgebox}

\subsection{稳定扩展为什么重要}

规模化模型的优势不仅是“参数更多”。更关键的是，当数据和计算预算增加时，训练行为是否可预测，失败能否定位，正则化是否成熟，以及 loss 是否能在无需复杂 bootstrap 的情况下稳定下降。RL 的困难往往不是表达能力不足，而是目标、数据分布和 policy 改变彼此耦合，使训练系统持续追逐自己生成的新分布。

\lecturefigure{slide-03-scalable-models-stable-training.jpg}{大规模 sequence model 的训练曲线通常更平滑，工程上更容易扩展。}{Stanford Online 官方视频 02:20；slide 引用 Kaplan et al. 2020。}

\teachervoice{讲者强调，stable training dynamics 对 practitioner 很重要：只要有更多数据和计算，就能逐步扩大模型并观察相对平滑的 loss 改善。Decision Transformer 的动机之一，是把这种“可训练性”带入 decision making，而不是只追求一个新的网络模块。}

\begin{warningbox}{稳定 loss 不等于稳定决策质量}
监督损失下降只能说明模型更好地拟合训练分布中的动作条件分布。真实环境回报还取决于 rollout 产生的状态分布、错误累积、目标回报是否可达以及环境随机性。后文必须单独测量 realized environmental return，不能用训练 loss 代替 policy evaluation。
\end{warningbox}

\subsection{从感知到行动}

感知模型输出标签、文本或表示；决策模型的输出会改变环境，并决定下一步将看到什么。这个闭环让 sequential decision making 比静态预测多出 distribution shift 和 credit assignment。讲者用“装有视觉和语言知识的圆柱体”指向机器人，不是为了说机器人只差一个 Transformer，而是为了强调 perception knowledge 必须最终转成 action sequence。

\lecturefigure{slide-04-sequential-data-to-sequential-decision.jpg}{研究问题从 sequential data 进一步转向 sequential decision making。}{Stanford Online 官方视频 03:10。}

\begin{importantbox}{核心转向}
静态序列建模学习 $p(x_t\mid x_{<t})$；决策序列建模需要学习
\[
p(a_t\mid \text{goal},s_{\le t},a_{<t},r_{<t}),
\]
其中动作 $a_t$ 会改变未来状态 $s_{t+1}$。因此“tokenization 相似”并不消除 closed-loop control 的风险，只是提供了一种统一的函数逼近和训练接口。
\end{importantbox}

\subsection{本章小结}

Decision Transformer 的出发点是训练范式迁移：把 RL 轨迹转成序列，希望复用 sequence modeling 的稳定目标与规模化能力。但决策具有闭环反馈，不能仅凭静态领域的成功证明方法有效，必须进入 MDP、离线数据和真实 rollout 评价。

